\documentclass{article}
\usepackage[utf8]{inputenc}
\usepackage[T1]{fontenc}
\usepackage{geometry}
\geometry{
    a4paper,
    left=2.5cm,
    right=2.5cm,
    top=2.5cm,
    bottom=2.5cm
}
\usepackage{graphicx}
\usepackage{float}
\usepackage{hyperref}
\usepackage[english]{babel}
\usepackage{amsmath}
\usepackage{amssymb}
\usepackage{array}
\usepackage{booktabs}
\usepackage{caption}
\usepackage{subcaption}
\usepackage{enumitem}
\usepackage{fancyhdr}
\usepackage{titlesec}
\usepackage{url}
\usepackage{textcomp}
\usepackage{xcolor}
\usepackage{listings}
\usepackage{multirow}
\usepackage{longtable}
\usepackage{tabularx}
\usepackage{times}

% Set up page style
\pagestyle{fancy}
\fancyhf{}
\rfoot{\thepage}
\renewcommand{\headrulewidth}{0pt}

% Title formatting
\titleformat{\chapter}[display]{\normalfont\bfseries}{\chaptertitle\thechapter}{0pt}{\Huge}
\titlespacing*{\chapter}{0pt}{50pt}{40pt}

% Document info
\title{Template Document}
\author{Academic Support Agent}
\date{\today}

\begin{document}
\maketitle

\pdfbookmark[0]{Title page}{title} % Sets a PDF bookmark for the title page
\setlength{\parindent}{0pt} % set paragraph indent
\setlength{\parskip}{0ex} % set paragraph spacing
\pagenumbering{alph} % prevents warnings about duplicated page numbering
% \renewcommand{\@pnumwidth}{1.9em} % was 1.55em
% \renewcommand{\@tocrmarg}{2.9em plus1fil} % raggedright toc entries (no hyphenation)

% Content from FrontMatter/FrontMatter would be included here


%%%%%%%%%%%%%%%%% MAIN %%%%%%%%%%%%%%%%%


% change header
\renewcommand{\chaptermark}[1]{\markboth{\chaptertitlename\ \thechapter: #1}{}}

\pagenumbering{arabic}
\setlength\abovedisplayshortskip{0.5\baselineskip}
\setlength\belowdisplayshortskip{0.5\baselineskip}
\setlength\abovedisplayskip{0.5\baselineskip}
\setlength\belowdisplayskip{0.5\baselineskip}

%%%%%%%%%%%%%%%%%%%%%
%%%%%%%%%%%%%%%%%%%%%
% For example

%% Content from Introduction would be included here


% The argument for \include is a filename of a LaTeX file contained within the same folder at this file. It could also be placed in a subfolder; the argument can then be specified as {subfolder_name/filename}.

% The \include'd file contains LaTeX commands that will, more or less, be read as if the code was contained in this file at the point the \include was called.

% There is an equivalent command, \includeonly, which you should place at the top of your document. Use this command to choose the files to include (rather than commenting lines). In this way, the document will be typeset as if the files were present (so links and cross-references won't be broken), but any output associated with the files will be excluded.

% Another tip: If using TeXShop, which uses synctex, you can command click on the preview to find the corresponding commands in your source. The reverse operation is also permitted. To do this from \include'd files, you need to tell TeXShop what the root file is. This is done by placing this line AT THE TOP of each \include'd file (you must include the "%"):
% !TEX root = root_filename.tex
% Doing this has other advantages, like being able to typeset from an included file, rather than having to typeset from the root file.

% Chapter 1
% Content from Chapter1/Chapter1 would be included here


% Chapter 2
% Content from Chapter2/Chapter2 would be included here


%% Appendices
%% appendices work in an identical way to chapters (in terms of syntax), but are labelled alphabetically rather than numerically

% Replace "Chapter" with "Appendix"
\addtocontents{toc}{\protect\renewcommand{\protect\cftchappresnum}{
\addtocontents{toc}{\protect\addtolength{\cftchapnumwidth}{1em}}
%% Content from Appendix1 would be included here

%% Content from Appendix2 would be included here


% Last I checked, the Appendix and Titlesec package do not get on with each other. So use "input" for appendix chapters instead. If you still want to use \includeonly, put the \input commands in another file and \include that.

%%%%%%%%%%%%%%%%% BACK MATTER %%%%%%%%%%%%%%%%%


% header
\renewcommand{\chaptermark}[1]{\markboth{#1}{}}

% Content from BackMatter/BackMatter would be included here
\end{document}